\section{Graph models}
\label{ch:graph}

The graph models read a dynamic graph between the wallets and the markets they trade, by
message passing at every block.  It is the
clearest negative result of the thesis, and this section gives the result at every stage of
the design, the two direct tests of what wallets know, the market--market view, and the
reasons.

\subsection{The wallet--market graph at every stage}
\label{sec:graph:stages}
\label{sec:exp:graph}


\Cref{tab:exp:graph} isolates the contribution of the wallet--market edges at every
stage.  It is never above $+0.3$ at 1h, it changes sign across configurations at 12h,
and in the best configuration it is $-0.1$ [$-0.2$, $0.0$] at 1h and $-1.0$ [$-1.2$,
$-0.8$] at 12h; at 6h it is $+0.65$ [$+0.20$, $+1.11$] on all samples but $-0.54$
[$-0.83$, $-0.27$] on the deduplicated subset.  The edge-orientation fix moves the
graph effect on the affected (token-2) markets from $-0.6/-0.3$ to $+0.05/+0.4$ points
at 1h/12h in a paired short test, with no overall gain because those markets are 4--9\%
of the samples.  Wallet-skill features never beat an activity-matched placebo
($-0.4$ at 1h, $0.0$ at 12h).  
\begin{table}[t]
\centering
\caption{Graph contribution at each stage (LOB-only / graph, all samples, seed 0).}
\label{tab:exp:graph}
\begin{fittable}\begin{tabular}{lcccc}
\toprule
Configuration & 1h & $\Delta$ & 12h & $\Delta$ \\
\midrule
lookback                      & 0.5038 / 0.5037 & $-0.01$ & 0.4368 / 0.4400 & $+0.32$ \\
+ hierarchical memory         & 0.5161 / 0.5184 & $+0.23$ & 0.4706 / 0.4559 & $-1.47$ \\
+ direct path + global state  & 0.5452 / 0.5461 & $+0.09$ & 0.4727 / 0.4724 & $-0.03$ \\
+ five lookbacks (best)       & 0.5547 / 0.5536 & $-0.11$ [$-0.21$, $0.00$] & 0.4914 / 0.4814 & $-0.99$ [$-1.15$, $-0.81$] \\
\bottomrule
\end{tabular}\end{fittable}
\end{table}

\subsection{Smart money}
\label{sec:graph:smartmoney}

A direct test of ``smart money'' outside the model confirms the picture.  For every wallet
the past hit rate at the horizon, the share of its trades on the side the price then moved
to, is computed on the training days and compared with the validation day: the Spearman
correlation is $0.5$--$0.6$, so skill is strongly persistent at the wallet level
(\cref{sec:expl:wallets}).  As a market-level predictor, the activity-weighted skill of the
wallets that traded a market in the last hour adds only about one point over an
identity-free placebo (random wallets with the same activity) and does not beat Mean
Reversion.  Inside the model, the same record given as an edge feature never beats the
placebo skill ($-0.4$ at 1h, $0.0$ at 12h), and the wallet's past behaviour (activity, maker
share, one-sidedness, size, breadth) gives $\pm 0.3$ with no consistent sign.

\subsection{The market--market view: sibling contracts}
\label{sec:graph:siblings}

The ten sibling features of \cref{sec:prep:features} were added to the control's table.  They lose $-0.10$ [$-0.18$,
$-0.02$] on all samples and $-0.25$ on the political markets, tie on the sports markets of the
same game ($-0.06$ [$-0.28$, +0.15]), and the best of them, the sum of the siblings' mids, is the
twenty-fifth feature by gain.  At one hour, after the market's own returns and activity clock,
what its siblings did in the last hour does not change the decision.  This is also the answer to
whether a market-to-market graph would help the deep model at this horizon: the information
such edges would carry is the information these features carry, and the trees found none; and
the lead--lag statistic of \cref{sec:expl:siblings} shows that at the minute scale, where the
siblings of a game might move first, they do not: the common cause reaches both books within
the same minute.

\subsection{Why the graph does not help}
\label{sec:graph:why}

The bipartite transaction graph is the clearest negative result of the thesis.  Four
observations explain it.
\begin{enumerate}[nosep]
  \item \textbf{The temporal signal dominates.}  Every improvement in
  \cref{sec:stateful:chain} came from temporal components, and once they are strong the
  graph's small positive effects at 12h in weaker configurations disappear.
  \item \textbf{Wallet memory is short.}  The decay is shared by wallets and markets, so
  a wallet state is erased within a minute of its last trade; most wallets trade far
  less often than that, and their state carries no history when they reappear.  A
  separate slow decay for wallets was proposed and not tested; the smart-money test
  suggests it would not change the conclusion.
  \item \textbf{Wallet skill is real but not market-level information.}  Past accuracy
  persists strongly at the wallet level, but a market is traded by many wallets and
  their aggregate is close to the placebo of random wallets with the same activity.
  \item \textbf{Most of the signal a wallet brings is already in the book.}  A trade is
  a book event; the order-flow features see it before the edge does.
\end{enumerate}
In sum, the temporal component weighs far more than the graph on this data.  The same conclusion emerges
from the literature on spatiotemporal benchmarks \cite{martin2026audit}: per-node
baselines are hard to beat and a graph must prove its value against them.

What would have to be different for a graph to help is also clear from the data.  A wallet
memory at a time scale of days rather than seconds, so that a wallet's state survives
between its trades; an edge feature that is not already a book event; or a market--market
graph on a question where the common cause reaches one book before the other, which the
lead--lag statistic did not find at the minute scale in this data.  A learned, sparse
market--market graph restricted to the markets of the same event, in the spirit of
\textcite{cini2023sparse}, remains the version of the idea worth testing.

